\documentclass[11pt,a4paper]{article}

\usepackage[T1]{fontenc}
\usepackage[utf8]{inputenc}
\usepackage[margin=2.5cm]{geometry}
\usepackage{amsmath,amssymb}
\usepackage{booktabs}
\usepackage{multirow}
\usepackage{xcolor}
\usepackage{hyperref}
\usepackage{parskip}
\usepackage{enumitem}
\setlist{nosep, topsep=3pt}

\hypersetup{colorlinks=true, urlcolor=blue!70!black, citecolor=blue!70!black}

\newcommand{\anonA}{[Anonymous,~2026a]}
\newcommand{\anonB}{[Anonymous,~2026b]}
\newcommand{\anonC}{[Anonymous,~2015]}
\newcommand{\anonD}{[Anonymous,~2017]}
\newcommand{\anonE}{[Anonymous,~2024]}

\begin{document}

% ---------------------------------------------------------------
\begin{center}
  {\LARGE\bfseries Gate-Based Quantum Neighborhood Search\\[4pt]
  for the Permutation Flow Shop Problem\\[6pt]
  on the Willow Processor}\\[12pt]
  {\large Willow Early Access Program --- Research Proposal}\\[4pt]
  {\normalsize Anonymous Submission}
\end{center}

\bigskip\hrule\bigskip

% ---------------------------------------------------------------
\section*{Scientific Motivation}

The Permutation Flow Shop Scheduling Problem (PFSP) is a classical NP-hard
combinatorial problem in which $n$ jobs must be processed on $m$ machines
in identical order, minimizing makespan $C_{\max}$. Neighborhood-based
local search is the dominant practical approach for large instances.

\paragraph{Prior work on classical neighborhoods.}
In \anonA{} we introduced and analyzed four neighborhood structures of
increasing complexity for the PFSP: \emph{Adjacent} (single adjacent
swap), \emph{Fibonacci} (optimal set of non-overlapping adjacent swaps
selected via dynamic programming, with solution count growing as Fibonacci
numbers), \emph{Dynasearch} (optimal set of non-overlapping endpoint swaps
over arbitrary segments), and \emph{Motzkin} (optimal set of non-crossing
nested endpoint swaps, with solution count given by Motzkin numbers $M_n$).
Experiments on Taillard benchmark instances with $n \in \{50,100,200,500\}$
in Iterated Local Search and Simulated Annealing showed that richer
neighborhoods (Dynasearch, Motzkin) consistently achieve lower makespan
within identical time budgets, at the cost of higher per-iteration
complexity.
The Head+Tail makespan evaluation technique employed in all four
neighborhoods builds on the block-decomposition framework for flow
shop acceleration developed in~\anonC{} and applied to parallel
metaheuristics in~\anonD{}.

\paragraph{Prior work on quantum neighborhoods.}
In \anonB{} we formulated all four neighborhoods as QUBO problems and
executed them on D-Wave quantum annealing hardware, providing the first
empirical comparison of classical and hybrid quantum variants under
identical wall-clock time budgets on Taillard instances ($n=20$, $m=5$).
The results revealed a large and consistent performance gap: the best
quantum variant achieved average RPD (Relative Percentage Deviation from
best-known makespans) of $18.4\%$, versus $2.6\%$ for classical Dynasearch
at a 10-second budget. The gap was traced to D-Wave hardware
overhead---cloud network latency, QPU queue time, and embedding
cost---which drastically reduces the number of neighborhood evaluations
per second relative to the CPU baseline.
This contrasts with the hybrid quantum-classical branch-and-bound
approach of~\anonE{}, which achieves competitive results for single-machine
scheduling by tightly coupling D-Wave with classical bounding, suggesting
that problem structure---not quantum hardware per se---determines
the achievable quality.

\paragraph{This proposal.}
The results of \anonB{} leave open a critical question:

\begin{center}
\emph{Does the classical--quantum performance gap arise from fundamental
limits of quantum optimization, or from the overhead of the annealing
paradigm on cloud hardware?}
\end{center}

The Willow processor is uniquely suited to answer this. Gate-based QAOA
execution eliminates cloud and embedding overhead entirely. However, unlike
quantum annealing, gate-based circuits accumulate error with each two-qubit
gate. This imposes a hard constraint on circuit depth that differs
fundamentally across the four neighborhood structures, depending on the
density of the corresponding QUBO interaction graph. We performed a
detailed feasibility analysis (error rate 0.14\% per CZ gate, Device~2
specification) to determine the tractable experimental regime for each
neighborhood.

% ---------------------------------------------------------------
\section*{Feasibility Analysis}

The number of CZ gates per QAOA layer equals twice the number of nonzero
off-diagonal entries $|J_{ij}|$ in the Ising Hamiltonian $H_C$. The
accumulated circuit error at depth $p$ is:
\[
\varepsilon(p) = 1 - (1 - \varepsilon_{\mathrm{CZ}})^{2\,|J_{ij}|\,p},
\qquad \varepsilon_{\mathrm{CZ}} = 0.0014.
\]
Table~\ref{tab:feasibility} summarizes the analysis across all four
neighborhoods and instance sizes. We label a configuration \emph{feasible}
if $\varepsilon < 10\%$, \emph{marginal} if $10\% \le \varepsilon < 25\%$
(used with ZNE mitigation), and \emph{infeasible} otherwise.

\begin{table}[h]
\centering
\caption{Accumulated CZ error per circuit, feasibility classification,
and assigned experimental role. $K$ = number of QUBO variables (qubits).
CZ/layer computed from QUBO graph structure (see text). ZNE = Zero-Noise
Extrapolation applied for marginal configurations.}
\label{tab:feasibility}
\smallskip
\begin{tabular}{@{}llrrrrll@{}}
\toprule
Neighborhood & $n$ & $K$ & CZ/layer & $p$ & $\varepsilon(p)$ &
  Status & Role \\
\midrule
Fibonacci  & 10 & 9  & 16  & 1 &  2.2\% & feasible   & primary \\
Fibonacci  & 10 & 9  & 16  & 2 &  4.4\% & feasible   & primary \\
Fibonacci  & 10 & 9  & 16  & 3 &  6.5\% & feasible   & primary \\
Fibonacci  & 20 & 19 & 36  & 1 &  4.9\% & feasible   & primary \\
Fibonacci  & 20 & 19 & 36  & 2 &  9.6\% & feasible   & primary \\
Fibonacci  & 30 & 29 & 56  & 1 &  7.5\% & feasible   & scaling \\
Fibonacci  & 50 & 49 & 96  & 1 & 12.6\% & marginal+ZNE & scaling \\
Fibonacci  & 100& 99 & 196 & 1 & 24.0\% & marginal+ZNE & scaling \\
\midrule
Adjacent   & 10 & 9  & 72  & 1 &  9.6\% & feasible   & baseline \\
\midrule
Motzkin    & 6  & 15 & 42  & 1 &  5.7\% & feasible   & small-$n$ \\
Motzkin    & 7  & 21 & 84  & 1 & 11.1\% & marginal+ZNE & small-$n$ \\
\midrule
Dynasearch & 5  & 10 & 22  & 1 &  3.0\% & feasible   & small-$n$ \\
Dynasearch & 6  & 15 & 52  & 1 &  7.0\% & feasible   & small-$n$ \\
\midrule
Fibonacci  & 20 & 19 & 36  & 3 & 14.0\% & marginal   & (excluded) \\
Adjacent   & 20 & 19 & 342 & 1 & 38.1\% & infeasible & excluded \\
Dynasearch & 10 & 45 & $\gg$1000 & 1 & $\approx$100\% & infeasible & excluded \\
Motzkin    & 10 & 45 & $\gg$1000 & 1 & $\approx$100\% & infeasible & excluded \\
\bottomrule
\end{tabular}
\end{table}

The analysis reveals a fundamental asymmetry: \emph{Fibonacci} is the
only neighborhood that scales to large $n$ on current gate-based hardware,
owing to its tridiagonal QUBO structure (only $n-2$ nonzero off-diagonal
entries). \emph{Adjacent} is feasible only at $n=10$ due to its dense,
all-to-all penalty structure. \emph{Dynasearch} and \emph{Motzkin}, with
$K=O(n^2)$ variables and dense conflict graphs, are restricted to
$n \le 6$--$7$ on Willow. This boundary is itself a publishable result.

% ---------------------------------------------------------------
\section*{Experiment Design}

\subsection*{Instance structure}

The experiment is organized in three groups reflecting the feasibility
analysis above.

\paragraph{Group A --- Fibonacci, scaling study ($n = 10, 20, 30, 50, 100$).}
Fibonacci is the primary experimental subject. Its tridiagonal QUBO
maps naturally to a nearest-neighbor chain on the Willow qubit grid,
minimizing routing overhead. Instances for $n \in \{10, 20, 30\}$ use
standard Taillard benchmarks generated with the published Taillard
generator (5 machines); instances for $n \in \{50, 100\}$ are generated
with the same generator. Circuit depths $p \in \{1,2,3\}$ for $n \le 20$;
$p = 1$ for $n \ge 30$ (feasibility constraint). At $n=100$, $K=99$
qubits, the QAOA simulation requires $2^{99}$ amplitudes classically ---
genuinely beyond the simulable regime, making Willow results irreplaceable.

\paragraph{Group B --- Adjacent, baseline ($n = 10$, $p = 1$).}
Adjacent at $n=10$ provides a direct comparison point against D-Wave
results from \anonB{} and classical ILS/SA results from \anonA{} at the
same instance size. It is included as a baseline, not scaled to larger $n$
due to its dense QUBO graph ($\varepsilon = 38\%$ already at $n=20$,
$p=1$).

\paragraph{Group C --- Dynasearch and Motzkin, small-$n$ ($n \le 7$).}
To include all four neighborhoods in the gate-based comparison, Dynasearch
($n \le 6$) and Motzkin ($n \le 7$) are executed on small instances.
Since no standard Taillard benchmarks exist for $n < 10$, instances are
generated using the same Taillard random generator with $m = 5$ machines
and seeds fixed for reproducibility. This group establishes that
gate-based QAOA can solve all four QUBO neighborhood formulations, and
provides the first comparison of the four neighborhoods under a gate-based
quantum model, even if at small scale.

\subsection*{QUBO and Ising formulation}

Each neighborhood defines a cost matrix $Q$ over $K$ binary variables
(full derivations in \anonB{}):
\[
\min_{\mathbf{x}\in\{0,1\}^K}\;\mathbf{x}^\top Q\,\mathbf{x}
\;\xrightarrow{\;x_i = \frac{1-Z_i}{2}\;}\;
H_C = \sum_i h_i Z_i + \sum_{i<j} J_{ij} Z_i Z_j.
\]
Linear terms $h_i$ encode swap costs $\delta_i$ computed via the
Head+Tail makespan technique; quadratic terms $J_{ij}$ encode conflict
penalties. The mapping is exact.

\medskip
\noindent\textbf{Adjacent} ($K=n-1$, dense $Q$): $Q_{ii} = \delta_i - P$,
$Q_{ij} = 2P$ for all $i \neq j$, with $P = 2\max_i|\delta_i|+1$.

\noindent\textbf{Fibonacci} ($K=n-1$, tridiagonal $Q$): $Q_{ii} = \delta_i$,
$Q_{i,i+1} = P$, all other entries zero. Only $n-2$ nonzero off-diagonal
entries.

\noindent\textbf{Dynasearch} ($K=O(n^2)$): variables for all segment
endpoint-swap pairs; $Q$ penalizes overlapping segments.

\noindent\textbf{Motzkin} ($K=O(n^2)$): variables for all endpoint-swap
pairs; $Q$ penalizes crossing and shared-endpoint pairs; nested and
disjoint pairs unconstrained.

\subsection*{QAOA circuits}

We implement QAOA of depth $p$:
\[
|\psi(\boldsymbol{\gamma},\boldsymbol{\beta})\rangle =
\prod_{l=1}^{p}
e^{-i\beta_l \sum_i X_i}\,
e^{-i\gamma_l H_C}\,
|{+}\rangle^{\otimes K}.
\]
Gate decomposition uses only native Willow gates: each $J_{ij}Z_iZ_j$
term maps to CZ\,--\,$R_Z(2\gamma J_{ij})$\,--\,CZ (two CZ gates); each
$h_i Z_i$ term maps to $R_Z(2\gamma h_i)$; the mixer maps to $R_X(2\beta)$
per qubit. No mid-circuit measurements, no adaptive feedback, no error
correction, no two-qubit gates other than CZ.

Variational parameters $(\boldsymbol\gamma, \boldsymbol\beta)$ are
optimized classically via gradient-free COBYLA with multi-start
initialization pre-warmed from analytic $p=1$ solutions.

% ---------------------------------------------------------------
\section*{Hardware Feasibility Summary}

All circuits in Groups A--C fit within the ${\sim}105$-qubit Willow grid.
The deepest feasible circuit (Fibonacci $n=20$, $p=2$: $K=19$, 72 CZ
gates total) runs in ${\sim}0.25\,\mu$s, a factor of ${\sim}400\times$
below the T1\,=\,98\,$\mu$s coherence limit (Device~2). Marginal
configurations ($n \ge 50$, $p=1$) reach at most ${\sim}0.7\,\mu$s.
Fibonacci at any $n$ maps to a linear qubit chain, requiring minimal
routing on the 2D Willow grid. Adjacent at $n=10$ requires all-to-all
connectivity within 9 qubits, achievable via SWAP routing within the
Willow grid at modest overhead.

\paragraph{Throughput.}
Total distinct circuits: Group~A ($5$ sizes $\times$ $\{1\text{--}3\}$
depths $\times$ 3 ZNE scales $\times$ 10 seeds $= 300$--$450$) $+$
Group~B ($1\times1\times3\times10 = 30$) $+$ Group~C ($4$ configs
$\times3\times10 = 120$). Total $\approx600$ circuits. At 60 circuits/s:
$\sim$10\,s QPU time. Well within one day.

% ---------------------------------------------------------------
\section*{Noise Mitigation}

We apply Zero-Noise Extrapolation (ZNE) via gate folding at
$\lambda \in \{1, 2, 3\}$, with Richardson polynomial extrapolation
to $\lambda=0$. The noise-impact observable
$\Delta_{\text{ZNE}} = \langle H_C\rangle_{\text{raw}} -
\langle H_C\rangle_{\text{ZNE}}$
is reported for all configurations and plotted as a function of $K$ and
$p$, providing empirical characterization of noise scaling across the
four neighborhood structures.

% ---------------------------------------------------------------
\section*{Observables and Planned Figures}

All observables are derived from bit-string measurement outcomes.

\begin{enumerate}

  \item \textbf{RPD vs.\ paradigm (primary result).}
        Average $\mathrm{RPD} = (C_{\max} - C^*)/C^* \times 100\%$
        for Classical CPU \anonA{}, D-Wave \anonB{}, and Willow QAOA,
        compared at identical $n$ and neighborhood. For Groups A and B
        ($n \le 20$, Taillard instances), this is a direct three-way
        comparison under identical QUBO formulations. This figure
        answers the central question of the proposal.

  \item \textbf{RPD vs.\ instance size --- Fibonacci scaling.}
        RPD as a function of $n \in \{10, 20, 30, 50, 100\}$ for
        Fibonacci QAOA versus classical ILS baseline from \anonA{}.
        At $n \ge 50$, the Willow results are genuinely beyond the
        classically simulable regime ($2^{49}$--$2^{99}$ amplitudes).
        This is the novel hardware-only result.

  \item \textbf{Energy vs.\ circuit depth.}
        $\langle H_C\rangle$ as a function of $p \in \{1,2,3\}$
        with and without ZNE, for Fibonacci at $n=10$ and $n=20$.

  \item \textbf{Noise impact vs.\ qubit count.}
        $\Delta_{\text{ZNE}}$ as a function of $K$ across all groups,
        quantifying empirical noise scaling with problem size.

  \item \textbf{All-neighborhood comparison at small $n$.}
        RPD for all four neighborhoods (Groups A--C) at the largest
        mutually feasible instance size, establishing a gate-based
        baseline across the full neighborhood hierarchy from \anonA{}
        and \anonB{}.

\end{enumerate}

% ---------------------------------------------------------------
\section*{Expected Scientific Impact}

\paragraph{Answering the overhead question.}
If QAOA on Willow achieves lower RPD than D-Wave for identical QUBO
instances, this is direct empirical evidence that the performance gap
reported in \anonB{} is driven by hardware overhead rather than
algorithmic limitations of quantum optimization. This would be a
significant and concrete result.

\paragraph{Novel scaling result.}
The Fibonacci scaling study (Group~A) addresses a question inaccessible
to both classical simulation and D-Wave: how does gate-based quantum
neighborhood search scale with $n$ for a linearly-scaling QUBO
neighborhood? At $n=100$, $K=99$ qubits, the Willow results are
irreplaceable by any classical means.

\paragraph{Structural insight.}
The contrast between Adjacent (dense $Q$) and Fibonacci (tridiagonal $Q$)
at $n=10$, $K=9$ isolates the effect of QUBO interaction graph sparsity
on QAOA performance --- a theoretically important question answered on
hardware for a practically motivated problem.

\paragraph{Hardware boundary.}
The feasibility analysis precisely maps where each of the four
neighborhoods becomes intractable on current gate-based hardware. This
boundary is itself a publishable result and a concrete benchmark for
future devices.

\paragraph{Publication target.}
\textit{npj Quantum Information}, \textit{Physical Review Applied},
or a top operations research venue.

% ---------------------------------------------------------------
\section*{Dedicated Researcher}

A dedicated PhD-level researcher will be exclusively assigned to this
experiment. The researcher has implemented all four QUBO neighborhood
formulations from scratch and validated them experimentally on D-Wave
hardware (\anonB{}). The full classical pipeline is implemented and
tested: QUBO matrix construction, $C_{\max}$ evaluation via the Head+Tail
technique, RPD benchmarking against Taillard instances, and a random
instance generator for small-$n$ experiments. The implementation of QAOA
circuits in Cirq, noise mitigation via ZNE, and integration with the
Willow processor constitute the planned contribution within this project.

A prototype QAOA implementation for the Fibonacci neighborhood
has been developed and validated on the Cirq simulator, including
circuit construction, COBYLA parameter optimization, and
end-to-end makespan improvement on Taillard-generated instances.
This codebase serves as the starting point for Willow integration.

% ---------------------------------------------------------------
\bigskip\hrule\bigskip

\paragraph{Note on prior work.}
\anonA{} and \anonB{} are two accepted conference papers (ICAISC 2026 and SOCO 2026)
from the proposing research group, withheld here to preserve anonymity.
Both are available upon request.
\anonC{} and \anonD{} are peer-reviewed journal publications in
\textit{Computers \& Industrial Engineering} and
\textit{Computers \& Operations Research} by the same group,
describing the block-decomposition acceleration framework and
parallel metaheuristic implementations underlying the present work.
\anonE{} is a recent publication in \textit{Future Generation Computer
Systems} from the same group on quantum annealing for scheduling,
available upon request.

\end{document}
